\documentclass[12pt, a4paper]{ctexart}
\usepackage{fontspec}
\usepackage{xcolor}
\usepackage{hyperref}
\usepackage{geometry}
\usepackage{enumitem}
\usepackage{parskip}
\usepackage{titlesec}
\usepackage{tcolorbox}
\tcbuselibrary{breakable}
\tcbset{autoparskip}

\usepackage{setspace}
\usepackage{listings}
\usepackage{amssymb}

\geometry{left=2.5cm, right=2.5cm, top=2.5cm, bottom=2.5cm}

\setmainfont{Times New Roman}
\setCJKmainfont{SimSun}

\hypersetup{
    colorlinks=true,
    linkcolor=blue!70!black,
    urlcolor=cyan,
    pdftitle={Java API - 逐题精讲解义},
    pdfauthor={专升本-fifine老师}
}

\definecolor{myblue}{RGB}{30,100,200}
\definecolor{lightblue}{RGB}{230,240,255}
\definecolor{darkblue}{RGB}{0,50,120}

\newtcolorbox{bluebox}[1][]{
    breakable,
    colback=lightblue,
    colframe=myblue,
    coltitle=darkblue,
    fonttitle=\bfseries,
    title={#1},
    boxrule=0.8pt,
    arc=4pt,
    left=6pt,
    right=6pt,
    top=6pt,
    bottom=6pt,
    before skip=8pt,
    after skip=8pt
}

\usepackage{etoolbox}
\BeforeBeginEnvironment{bluebox}{\par\vfill}%
\AfterEndEnvironment{bluebox}{\par\vfill}%

\titleformat{\section}
  {\normalfont\Large\bfseries\color{myblue}}{\thesection}{1em}{}
\titlespacing*{\section}{0pt}{3.5ex plus 1ex minus .2ex}{1.5ex plus .2ex}

\title{\textcolor{myblue}{\textbf{Java API - 逐题精讲解义}}}
\author{专升本-fifine老师 教学组}
\date{2026年03月05日}

\lstset{
    language=Java,
    basicstyle=\ttfamily\small,
    keywordstyle=\color{blue},
    commentstyle=\color{green!60!black},
    stringstyle=\color{orange},
    numbers=left,
    numberstyle=\tiny\color{gray},
    frame=single,
    breaklines=true,
    columns=fullflexible,
    showstringspaces=false
}

\newcommand{\code}[1]{\texttt{#1}}

\begin{document}

\maketitle
\thispagestyle{empty}
\newpage

\section{Java API}

\begin{enumerate}[label=\textbf{\arabic*.}]

	\item \textbf{以下哪个方法可以将字符串转换为整数？}
	      \begin{enumerate}[label=\Alph*.]
		      \item parseInt()
		      \item valueOf()
		      \item toString()
		      \item toInteger()
	      \end{enumerate}

	      \begin{bluebox}{答案与深度解析}
		      \textbf{答案：A. parseInt()}

		      \textbf{【核心认知框架>}
		      \begin{itemize}[label={}, leftmargin=*, nosep]
			      \item \textbf{题目本质}：考查Integer类的静态方法功能区分
			      \item \textbf{关键区分点}：parseXxx返回基本类型，valueOf返回包装类对象
			      \item \textbf{命题人意图}：测试对包装类方法返回类型的理解
		      \end{itemize}

		      \textbf{【选项原理深度拆解】}

		      \begin{itemize}[leftmargin=*, nosep, label={}]
			      \item[A] \textbf{parseInt() — 正确（返回基本类型int）}
			            \begin{itemize}[label={}, leftmargin=*, nosep]
				            \item \textbf{字节码层面}：调用Integer.parseInt返回int类型，不产生对象
				            \item \textbf{方法签名}：public static int parseInt(String s) throws NumberFormatException
				            \item \textbf{使用场景}：需要基本类型计算时使用
				            \item \textbf{代码示例}：
				                  \begin{lstlisting}[language=Java]
int value = Integer.parseInt("123");
// 字节码：invokestatic Integer.parseInt
// 直接返回基本类型，无需拆箱
				            \end{lstlisting}
			            \end{itemize}

			      \item[B] \textbf{valueOf() — 错误（返回包装类Integer对象）}
			            \begin{itemize}[label={}, leftmargin=*, nosep]
				            \item \textbf{方法签名}：public static Integer valueOf(String s)
				            \item \textbf{返回类型}：Integer对象，不是int
				            \item \textbf{字节码}：调用Integer.valueOf后返回引用类型
				            \item \textbf{缓存机制}：-128~127范围内返回缓存对象
				            \item \textbf{代码示例}：
				                  \begin{lstlisting}[language=Java]
Integer obj = Integer.valueOf("123");
// 返回包装类对象，需用intValue()获取int值
				            \end{lstlisting}
			            \end{itemize}

			      \item[C] \textbf{toString() — 错误（反向转换）}
			            \begin{itemize}[label={}, leftmargin=*, nosep]
				            \item 功能：将其他类型转为字符串
				            \item 即Integer.toString(123)返回 String "123"
				            \item 不是字符串转整数
			            \end{itemize}

			      \item[D] \textbf{toInteger() — 错误（不存在此方法）}
			            \begin{itemize}[label={}, leftmargin=*, nosep]
				            \item Integer类没有toInteger()方法
				            \item 可能混淆了toString()和parseInt()
			            \end{itemize}
		      \end{itemize}

		      \textbf{【应背诵知识点>}
		      \begin{itemize}[label={}, nosep]
			      \item parseInt()：返回基本类型，无缓存
			      \item valueOf()：返回包装类，有缓存(-128~127)
			      \item 选项辨析：parse返回基本类型，valueOf返回对象，toString转字符串
		      \end{itemize}

		      \textbf{【命题趋势与实战策略】}

		      \textbf{考场秒杀}：看到"字符串转整数"→选parseInt()
		      \textbf{陷阱清单}：valueOf返回对象，需再调intValue()才得int
	      \end{bluebox}

	      \vspace{1em}

	\item \textbf{在Java中，以下哪个方法用于获取字符串的长度？}
	      \begin{enumerate}[label=\Alph*.]
		      \item length()
		      \item size()
		      \item count()
		      \item getLength()
	      \end{enumerate}

	      \begin{bluebox}{答案与深度解析}
		      \textbf{答案：A. length()}

		      \textbf{【核心认知框架>}
		      \begin{itemize}[label={}, leftmargin=*, nosep]
			      \item \textbf{题目本质}：区分字符串length()与集合size()
			      \item \textbf{关键区分点}：String用length，集合用size
			      \item \textbf{命题人意图}：测试API记忆的准确性
		      \end{itemize}

		      \textbf{【选项原理深度拆解】}

		      \begin{itemize}[leftmargin=*, nosep, label={}]
			      \item[A] \textbf{length() — 正确（字符串专用）}
			            \begin{itemize}[label={}, leftmargin=*, nosep]
				            \item \textbf{方法签名}：public int length()
				            \item \textbf{字节码}：底层调用String.value.length字段
				            \item \textbf{底层实现}：String类使用private final char[] value存储字符，length()返回value.length
				            \item \textbf{O(1)复杂度}：直接返回数组长度，不遍历
				            \item \textbf{代码示例}：
				                  \begin{lstlisting}[language=Java]
String s = "Hello";
int len = s.length();
// 字节码：getfield String.value, arraylength
// 只是字段访问，非常快
				            \end{lstlisting}
			            \end{itemize}

			      \item[B] \textbf{size() — 错误（集合方法）}
			            \begin{itemize}[label={}, leftmargin=*, nosep]
				            \item String类没有size()方法
				            \item size()是Collection接口方法
				            \item \textbf{常见混淆}：ArrayList.size()、HashSet.size()
			            \end{itemize}

			      \item[C] \textbf{count() — 错误（不存在）}
			            \begin{itemize}[label={}, leftmargin=*, nosep]
				            \item String类没有count()方法
				            \item 某些库有countChars()等，非Java标准
			            \end{itemize}

			      \item[D] \textbf{getLength() — 错误（不存在）}
			            \begin{itemize}[label={}, leftmargin=*, nosep]
				            \item 典型误记：把length写成了getLength()
			            \end{itemize}
		      \end{itemize}

		      \textbf{【应背诵知识点>}
		      \begin{itemize}[label={}, nosep]
			      \item String.length() → 返回字符数
			      \item Collection.size() → 返回元素个数
			      \item 数组.length → 属性，不是方法
		      \end{itemize}

		      \textbf{【命题趋势与实战策略】}

		      \textbf{考场秒杀}：看到"字符串长度"→选length()
		      \textbf{记忆口诀}：字符串长度length，集合大小size，数组长度length(属性)
	      \end{bluebox}

	      \vspace{1em}

	\item \textbf{以下哪个方法用于将字符串转换为大写？}
	      \begin{enumerate}[label=\Alph*.]
		      \item toUpperCase()
		      \item upperCase()
		      \item makeUpperCase()
		      \item convertToUpperCase()
	      \end{enumerate}

	      \begin{bluebox}{答案与深度解析}
		      \textbf{答案：A. toUpperCase()}

		      \textbf{【核心认知框架>}
		      \begin{itemize}[label={}, leftmargin=*, nosep]
			      \item \textbf{题目本质}：测试String类大小写转换方法
			      \item \textbf{关键区分点}：to前缀是Java标准命名
			      \item \textbf{命题人意图}：排除干扰项，记忆精确API
		      \end{itemize}

		      \textbf{【选项原理深度拆解】}

		      \begin{itemize}[leftmargin=*, nosep, label={}]
			      \item[A] \textbf{toUpperCase() — 正确（标准方法）}
			            \begin{itemize}[label={}, leftmargin=*, nosep]
				            \item \textbf{方法签名}：public String toUpperCase()
				            \item \textbf{返回新字符串}：原字符串不变（String不可变）
				            \item \textbf{底层实现}：使用Character.toUpperCase()遍历每个字符
				            \item \textbf{本地化支持}：可传Locale参数处理特殊语言
				            \item \textbf{代码示例}：
				                  \begin{lstlisting}[language=Java]
String s1 = "hello";
String s2 = s1.toUpperCase();
// s1仍为"hello"，s2为"HELLO"
// 因为String不可变，返回新对象

// 本地化版本
String s3 = "ı"; // 土耳其语小写i
s3.toUpperCase(); // "I" (错误)
s3.toUpperCase(Locale.forLanguageTag("tr")); // "İ" (正确)
				            \end{lstlisting}
			            \end{itemize}

			      \item[B] \textbf{upperCase() — 错误（不存在）}
			            \begin{itemize}[label={}, leftmargin=*, nosep]
				            \item 某些语言用upperCase()，Java使用to前缀
			            \end{itemize}

			      \item[C] \textbf{makeUpperCase() — 错误（不存在）}
			            \begin{itemize}[label={}, leftmargin=*, nosep]
				            \item GNU C++等库有类似命名，非Java
			            \end{itemize}

			      \item[D] \textbf{convertToUpperCase() — 错误（不存在）}
			            \begin{itemize}[label={}, leftmargin=*, nosep]
				            \item 过于冗长，不符合Java简洁命名习惯
			            \end{itemize}
		      \end{itemize}

		      \textbf{【应背诵知识点>}
		      \begin{itemize}[label={}, nosep]
			      \item toUpperCase()：转大写，返回新字符串
			      \item toLowerCase()：转小写，返回新字符串
			      \item String是不可变类，所有方法都返回新对象
		      \end{itemize}

		      \textbf{【命题趋势与实战策略】}

		      \textbf{考场秒杀}：看到转大写→选toUpperCase()
		      \textbf{记忆技巧}：Java字符串方法都用toXxx()模式
	      \end{bluebox}

	      \vspace{1em}

	\item \textbf{以下哪个方法用于将字符串转换为小写？}
	      \begin{enumerate}[label=\Alph*.]
		      \item toLowerCase()
		      \item lowerCase()
		      \item makeLowerCase()
		      \item convertToLowerCase()
	      \end{enumerate}

	      \begin{bluebox}{答案与深度解析}
		      \textbf{答案：A. toLowerCase()}

		      \textbf{【核心认知框架>}
		      \begin{itemize}[label={}, leftmargin=*, nosep]
			      \item \textbf{题目本质}：与toUpperCase对称，测试记忆对称性
			      \item \textbf{关键区分点}：to前缀 + 首字母大写的驼峰命名
			      \item \textbf{命题人意图}：确认对toLowerCase的准确记忆
		      \end{itemize}

		      \textbf{【选项原理深度拆解】}

		      \begin{itemize}[leftmargin=*, nosep, label={}]
			      \item[A] \textbf{toLowerCase() — 正确（标准方法）}
			            \begin{itemize}[label={}, leftmargin=*, nosep]
				            \item \textbf{方法签名}：public String toLowerCase()
				            \item \textbf{与toUpperCase对称}：唯一正确的转小写方法
				            \item \textbf{代码示例}：
				                  \begin{lstlisting}[language=Java]
String s = "HELLO";
String low = s.toLowerCase();
// low = "hello"，s仍为"HELLO"
				            \end{lstlisting}
			            \end{itemize}

			      \item[B/D] \textbf{其他选项 — 错误（不存在）}
			            \begin{itemize}[label={}, leftmargin=*, nosep]
				            \item 同上题，均为干扰选项
			            \end{itemize}
		      \end{itemize}

		      \textbf{【应背诵知识点>}
		      \begin{itemize}[label={}, nosep]
			      \item toUpperCase() / toLowerCase() 成对记忆
			      \item 都是to前缀 + 首字母大写 + Case()
		      \end{itemize}

		      \textbf{【命题趋势与实战策略】}

		      \textbf{考场秒杀}：转小写→toLowerCase()，与toUpperCase()对称
	      \end{bluebox}

	      \vspace{1em}

	\item \textbf{以下哪个方法用于截取字符串？}
	      \begin{enumerate}[label=\Alph*.]
		      \item substring()
		      \item cut()
		      \item slice()
		      \item truncate()
	      \end{enumerate}

	      \begin{bluebox}{答案与深度解析}
		      \textbf{答案：A. substring()}

		      \textbf{【核心认知框架>}
		      \begin{itemize}[label={}, leftmargin=*, nosep]
			      \item \textbf{题目本质}：考查字符串截取方法命名
			      \item \textbf{关键区分点}：substring暗示"子字符串"
			      \item \textbf{命题人意图}：区分JavaScript slice()与Java substring()
		      \end{itemize}

		      \textbf{【选项原理深度拆解】}

		      \begin{itemize}[leftmargin=*, nosep, label={}]
			      \item[A] \textbf{substring() — 正确（Java标准方法）}
			            \begin{itemize}[label={}, leftmargin=*, nosep]
				            \item \textbf{方法签名}：
				                  \begin{lstlisting}[language=Java]
public String substring(int beginIndex)
public String substring(int beginIndex, int endIndex)
				            \end{lstlisting}
				            \item \textbf{参数说明}：
				                  \begin{itemize}[label={}, nosep]
					                  \item 单参数：从beginIndex到末尾
					                  \item 双参数：[beginIndex, endIndex) 左闭右开
				                  \end{itemize}
				            \item \textbf{底层实现}：创建新String对象，共享原char[]（Java7之后复制）
				            \item \textbf{字节码}：调用String构造器
				            \item \textbf{代码示例}：
				                  \begin{lstlisting}[language=Java]
String s = "HelloWorld";
String s1 = s.substring(5);      // "World"
String s2 = s.substring(0, 5);   // "Hello"

// 边界检查
// s.substring(5, 5) → "" (空字符串)
// s.substring(0, 10) → "HelloWorld"
// s.substring(-1) → StringIndexOutOfBoundsException
// s.substring(11) → StringIndexOutOfBoundsException
				            \end{lstlisting}
				            \item \textbf{性能注意}：循环中频繁substring可能产生大量临时对象
			            \end{itemize}

			      \item[B] \textbf{cut() — 错误（不存在）}
			            \begin{itemize}[label={}, leftmargin=*, nosep]
				            \item 某些语言有cut()方法，Java没有
			            \end{itemize}

			      \item[C] \textbf{slice() — 错误（JavaScript方法）}
			            \begin{itemize}[label={}, leftmargin=*, nosep]
				            \item JavaScript Array/String有slice()方法
				            \item 跨语言陷阱：容易混淆
			            \end{itemize}

			      \item[D] \textbf{truncate() — 错误（不存在）}
			            \begin{itemize}[label={}, leftmargin=*, nosep]
				            \item truncate通常用于文件截断，非方法名
			            \end{itemize}
		      \end{itemize}

		      \textbf{【应背诵知识点>}
		      \begin{itemize}[label={}, nosep]
			      \item substring(begin)：从begin到末尾
			      \item substring(begin, end)：[begin, end) 左闭右开
			      \item 选项辨析：substring是字符串专用，slice是数组/JS方法
		      \end{itemize}

		      \textbf{【命题趋势与实战策略】}

		      \textbf{考场秒杀}：截取字符串→substring()
		      \textbf{高频陷阱}：endIndex是要截取到的位置左闭右开
	      \end{bluebox}

	      \vspace{1em}

	\item \textbf{以下哪个方法用于判断字符串是否以指定字符串开头？}
	      \begin{enumerate}[label=\Alph*.]
		      \item startsWith()
		      \item beginWith()
		      \item isStartWith()
		      \item hasStartWith()
	      \end{enumerate}

	      \begin{bluebox}{答案与深度解析}
		      \textbf{答案：A. startsWith()}

		      \textbf{【核心认知框架>}
		      \begin{itemize}[label={}, leftmargin=*, nosep]
			      \item \textbf{题目本质}：考查字符串前缀/后缀判断方法
			      \item \textbf{关键区分点}：Start/End + With + s（复数）
			      \item \textbf{命题人意图}：测试对方法命名习惯的理解
		      \end{itemize}

		      \textbf{【选项原理深度拆解】}

		      \begin{itemize}[leftmargin=*, nosep, label={}]
			      \item[A] \textbf{startsWith() — 正确（标准方法）}
			            \begin{itemize}[label={}, leftmargin=*, nosep]
				            \item \textbf{方法签名}：
				                  \begin{lstlisting}[language=Java]
public boolean startsWith(String prefix)
public boolean startsWith(String prefix, int toffset)
				            \end{lstlisting}
				            \item \textbf{底层实现}：比较String.value数组前缀字符
				            \item \textbf{性能}：O(k)复杂度，k为prefix长度
				            \item \textbf{代码示例}：
				                  \begin{lstlisting}[language=Java]
String filename = "document.pdf";
boolean isPdf = filename.startsWith(".pdf");
System.out.println(isPdf);  // false (注意位置!)

boolean isDoc = filename.startsWith("doc");
System.out.println(isDoc);  // true

// 带offset版本
String s = "HelloWorld";
boolean startsAt5 = s.startsWith("World", 5);  // true
boolean startsAt0 = s.startsWith("World", 0);  // false
				            \end{lstlisting}
				            \item \textbf{空字符串}：s.startsWith("") 永远返回true
				            \item \textbf{边界检查}：prefix为null抛出NPE
			            \end{itemize}

			      \item[B] \textbf{beginWith() — 错误（动词begin → start）}
			            \begin{itemize}[label={}, leftmargin=*, nosep]
				            \item verb naming: Start(begin) vs Begin
				            \item Java用start不用begin
			            \end{itemize}

			      \item[C] \textbf{isStartWith() — 错误（is前缀是布尔判断）}
			            \begin{itemize}[label={}, leftmargin=*, nosep]
				            \item Java中布尔方法一般用is开头（如isEmpty()）
				            \item 但startsWith()/endsWith()例外
			            \end{itemize}

			      \item[D] \textbf{hasStartWith() — 错误（语义不通）}
			            \begin{itemize}[label={}, leftmargin=*, nosep]
				            \item has通常用于"是否有"属性，不用于前缀判断
			            \end{itemize}
		      \end{itemize}

		      \textbf{【应背诵知识点>}
		      \begin{itemize}[label={}, nosep]
			      \item startsWith(prefix)：判断以前缀开始
			      \item endsWith(suffix)：判断以后缀结束（对称记忆）
			      \item 方法无s复数：StartWith是错误的
		      \end{itemize}

		      \textbf{【命题趋势与实战策略】}

		      \textbf{考场秒杀}：以前缀开头→startsWith()，与endsWith()成对
		      \textbf{记忆口诀}：Start开始End结束，都要加With+s
	      \end{bluebox}

	      \vspace{1em}

	\item \textbf{以下哪个方法用于判断字符串是否以指定字符串结尾？}
	      \begin{enumerate}[label=\Alph*.]
		      \item endsWith()
		      \item finishWith()
		      \item isEndWith()
		      \item hasEndWith()
	      \end{enumerate}

	      \begin{bluebox}{答案与深度解析}
		      \textbf{答案：A. endsWith()}

		      \textbf{【核心认知框架>}
		      \begin{itemize}[label={}, leftmargin=*, nosep]
			      \item \textbf{题目本质}：与startsWith对称，考查一致性记忆
			      \item \textbf{关键区分点}：End + With + s，与Start完整对应
			      \item \textbf{命题人意图}：测试方法的对称记忆能力
		      \end{itemize}

		      \textbf{【选项原理深度拆解】}

		      \begin{itemize}[leftmargin=*, nosep, label={}]
			      \item[A] \textbf{endsWith() — 正确（标准方法）}
			            \begin{itemize}[label={}, leftmargin=*, nosep]
				            \item \textbf{方法签名}：public boolean endsWith(String suffix)
				            \item \textbf{底层实现}：比较String.value数组后缀字符
				            \item \textbf{与startsWith对称}：都是End/Start With
				            \item \textbf{代码示例}：
				                  \begin{lstlisting}[language=Java]
String filename = "report.pdf";
boolean isPdf = filename.endsWith(".pdf");
System.out.println(isPdf);  // true

boolean isDoc = filename.endsWith(".doc");
System.out.println(isDoc);  // false

// 空字符串：始终true
"Hello".endsWith("");  // true
// null：抛出NPE
"Hello".endsWith(null);  // NullPointerException
				            \end{lstlisting}
			            \end{itemize}

			      \item[B/C/D] \textbf{其他选项 — 错误（不存在）}
			            \begin{itemize}[label={}, leftmargin=*, nosep]
				            \item 同startsWith题的干扰模式
			            \end{itemize}
		      \end{itemize}

		      \textbf{【应背诵知识点>}
		      \begin{itemize}[label={}, nosep]
			      \item endsWith(suffix)：判断以suffix结束
			      \item start/End + With + s 背诵公式
		      \end{itemize}

		      \textbf{【命题趋势与实战策略】}

		      \textbf{考场秒杀}：以后缀结束→endsWith()
		      \textbf{记忆技巧}：Start开始End结束，一对好基友一起记
	      \end{bluebox}

	      \vspace{1em}

	\item \textbf{以下哪个方法用于在字符串中查找指定字符第一次出现的位置？}
	      \begin{enumerate}[label=\Alph*.]
		      \item indexOf()
		      \item find()
		      \item search()
		      \item locate()
	      \end{enumerate}

	      \begin{bluebox}{答案与深度解析}
		      \textbf{答案：A. indexOf()}

		      \textbf{【核心认知框架>}
		      \begin{itemize}[label={}, leftmargin=*, nosep]
			      \item \textbf{题目本质}：考查字符串查找方法的命名
			      \item \textbf{关键区分点}：index（索引）vs search（搜索）
			      \item \textbf{命题人意图}：区分indexOf返回索引值，search返回布尔/对象
		      \end{itemize}

		      \textbf{【选项原理深度拆解】}

		      \begin{itemize}[leftmargin=*, nosep, label={}]
			      \item[A] \textbf{indexOf() — 正确（返回索引位置）}
			            \begin{itemize}[label={}, leftmargin=*, nosep]
				            \item \textbf{方法签名}：
				                  \begin{lstlisting}[language=Java]
public int indexOf(int ch)
public int indexOf(String str)
public int indexOf(int ch, int fromIndex)
public int indexOf(String str, int fromIndex)
				            \end{lstlisting}
				            \item \textbf{返回值规则}：
				                  \begin{itemize}[label={}, nosep]
					                  \item 找到：返回第一次出现的位置（从0开始）
					                  \item 未找到：返回-1（约定俗成）
				                  \end{itemize}
				            \item \textbf{底层实现}：遍历char[]数组，比较每个字符
				            \item \textbf{性能}：O(n)复杂度，n为字符串长度
				            \item \textbf{代码示例}：
				                  \begin{lstlisting}[language=Java]
String s = "HelloWorld";
System.out.println(s.indexOf('l'));      // 2 (第一个'l')
System.out.println(s.indexOf("World"));  // 5
System.out.println(s.indexOf('z'));      // -1 (不存在)

// 带fromIndex版本
String s = "ababab";
System.out.println(s.indexOf('b', 2));   // 3 (从索引2开始找)

// 注意：区分int和String参数
s.indexOf('H');  // 字符参数，找H的位置 → 0
s.indexOf("H");  // 字符串参数，找"H"的位置 → 0
				            \end{lstlisting}
				            \item \textbf{命题陷阱}：返回-1不是异常，0是有效索引
			            \end{itemize}

			      \item[B] \textbf{find() — 错误（不存在）}
			            \begin{itemize}[label={}, leftmargin=*, nosep]
				            \item C++的string::find，Java没有
			            \end{itemize}

			      \item[C] \textbf{search() — 错误（不存在）}
			            \begin{itemize}[label={}, leftmargin=*, nosep]
				            \item JavaScript有search()，返回匹配索引或-1
				            \item Java中Pattern/Matcher有find()返回boolean
			            \end{itemize}

			      \item[D] \textbf{locate() — 错误（不存在）}
			            \begin{itemize}[label={}, leftmargin=*, nosep]
				            \item locate通常用于定位文件、资源等
			            \end{itemize}
		      \end{itemize}

		      \textbf{【应背诵知识点>}
		      \begin{itemize}[label={}, nosep]
			      \item indexOf(char)：查找字符第一次出现，返回索引或-1
			      \item indexOf(String)：查找子串第一次出现
			      \item 未找到返回-1，找到返回第一次位置
		      \end{itemize}

		      \textbf{【命题趋势与实战策略】}

		      \textbf{考场秒杀}：查找位置→indexOf()
		      \textbf{高频陷阱}：返回-1不是0，0是第一个字符有效位置
		      \textbf{对仗记忆}：indexOf找第一次，lastIndexOf找最后一次
	      \end{bluebox}

	      \vspace{1em}

	\item \textbf{以下哪个方法用于在字符串中查找指定字符最后一次出现的位置？}
	      \begin{enumerate}[label=\Alph*.]
		      \item lastIndexOf()
		      \item findLast()
		      \item searchLast()
		      \item locateLast()
	      \end{enumerate}

	      \begin{bluebox}{答案与深度解析}
		      \textbf{答案：A. lastIndexOf()}

		      \textbf{【核心认知框架>}
		      \begin{itemize}[label={}, leftmargin=*, nosep]
			      \item \textbf{题目本质}：与indexOf对称，考查反方向查找
			      \item \textbf{关键区分点}：last + Index + Of，与indexOf对应
			      \item \textbf{命题人意图}：测试对正反向查找方法的完整记忆
		      \end{itemize}

		      \textbf{【选项原理深度拆解】}

		      \begin{itemize}[leftmargin=*, nosep, label={}]
			      \item[A] \textbf{lastIndexOf() — 正确（反向查找）}
			            \begin{itemize}[label={}, leftmargin=*, nosep]
				            \item \textbf{方法签名}：
				                  \begin{lstlisting}[language=Java]
public int lastIndexOf(int ch)
public int lastIndexOf(String str)
public int lastIndexOf(int ch, int fromIndex)
public int lastIndexOf(String str, int fromIndex)
				            \end{lstlisting}
				            \item \textbf{搜索方式}：从字符串末尾向前搜索
				            \item \textbf{返回值规则}：同indexOf（找到返回索引，未找到返回-1）
				            \item \textbf{底层实现}：逆序遍历char[]数组
				            \item \textbf{代码示例}：
				                  \begin{lstlisting}[language=Java]
String s = "ababab";
System.out.println(s.indexOf('a'));    // 0 (第一个)
System.out.println(s.lastIndexOf('a')); // 4 (最后一个)

// 查找子串
String s2 = "HelloHello";
System.out.println(s2.indexOf("Hello"));     // 0
System.out.println(s2.lastIndexOf("Hello")); // 5

// fromIndex参数：反向搜索的起点
String s3 = "banana";
// 从索引3开始向前搜索'a'
System.out.println(s3.lastIndexOf('a', 3)); // 1 (不是3)
// 解释：从索引3(即'n')向前找，遇到索引1的'a'
				            \end{lstlisting}
				            \item \textbf{常见场景}：解析路径、提取文件扩展名
			            \end{itemize}

			      \item[B/C/D] \textbf{其他选项 — 错误（不存在）}
			            \begin{itemize}[label={}, leftmargin=*, nosep]
				            \item 与indexOf类似，都是为干扰词
			            \end{itemize}
		      \end{itemize}

		      \textbf{【应背诵知识点>}
		      \begin{itemize}[label={}, nosep]
			      \item indexOf：从左往右找第一次
			      \item lastIndexOf：从右往左找最后一次
			      \item 都返回索引，找不到返回-1
			      \item 对仗记忆：第一次/最后一次 ↔ indexOf/lastIndexOf
		      \end{itemize}

		      \textbf{【命题趋势与实战策略】}

		      \textbf{考场秒杀}：最后一次出现→lastIndexOf()
		      \textbf{记忆口诀}：first用indexOf，last用lastIndexOf，找不到都是-1
	      \end{bluebox}

	      \vspace{1em}

	\item \textbf{以下关于Java中的包装类的说法错误的是}
	      \begin{enumerate}[label=\Alph*.]
		      \item 包装类可以将基本数据类型转换为对象
		      \item Integer是int的包装类
		      \item 包装类的对象不可变
		      \item 包装类可以提高基本数据类型的操作效率
	      \end{enumerate}

	      \begin{bluebox}{答案与深度解析}
		      \textbf{答案：D. 包装类可以提高基本数据类型的操作效率}

		      \textbf{【核心认知框架>}
		      \begin{itemize}[label={}, leftmargin=*, nosep]
			      \item \textbf{题目本质}：考查包装类的本质特性和适用场景
			      \item \textbf{关键区分点}：包装类带来的是"对象能力"而非"性能提升"
			      \item \textbf{命题人意图}：纠正"包装类=更好性能"的常见误解
		      \end{itemize}

		      \textbf{【选项原理深度拆解】}

		      \begin{itemize}[leftmargin=*, nosep, label={}]
			      \item[A] \textbf{包装类可以将基本数据类型转换为对象 — 正确}
			            \begin{itemize}[label={}, leftmargin=*, nosep]
				            \item \textbf{基本作用}：为基本类型提供对象形式
				            \item \textbf{应用场景}：
				                  \begin{itemize}[label={}, nosep]
					                  \item 泛型要求使用对象类型（List<Integer>不可用int）
					                  \item 集合只能存储对象
					                  \item 需要null值表示"无数据"
				                  \end{itemize}
				            \item \textbf{自动装箱/拆箱}：
				                  \begin{lstlisting}[language=Java]
Integer obj = 100;    // 自动装箱 int → Integer
int val = obj;        // 自动拆箱 Integer → int

// 编译器等价于：
Integer obj = Integer.valueOf(100);
int val = obj.intValue();

// 字节码验证：
// 自动装箱：invokestatic Integer.valueOf(I)Ljava/lang/Integer;
// 自动拆箱：invokevirtual Integer.intValue()I
				            \end{lstlisting}
			            \end{itemize}

			      \item[B] \textbf{Integer是int的包装类 — 正确}
			            \begin{itemize}[label={}, leftmargin=*, nosep]
				            \item \textbf{8大包装类}：
				                  \begin{tabular}{|l|l|}
					                  \hline
					                  \textbf{基本类型} & \textbf{包装类} \\ \hline
					                  byte          & Byte         \\ \hline
					                  short         & Short        \\ \hline
					                  int           & Integer      \\ \hline
					                  long          & Long         \\ \hline
					                  float         & Float        \\ \hline
					                  double        & Double       \\ \hline
					                  char          & Character    \\ \hline
					                  boolean       & Boolean      \\ \hline
				                  \end{tabular}
				            \item \textbf{命名单词差异}：int→Integer，char→Character，其他首字母大写
				            \item \textbf{原因}：避免与关键字冲突（Integer是完整拼写）
			            \end{itemize}

			      \item[C] \textbf{包装类的对象不可变 — 正确}
			            \begin{itemize}[label={}, leftmargin=*, nosep]
				            \item \textbf{不可变性证据}：
				                  \begin{lstlisting}[language=Java]
public final class Integer extends Number {
    private final int value;  // final字段

    public Integer(int value) {
        this.value = value;
    }
    // 没有setter方法
}
				            \end{lstlisting}
				            \item \textbf{字节码}：所有构造后value字段不可修改
				            \item \textbf{衍生结果}：
				                  \begin{itemize}[label={}, nosep]
					                  \item 修改操作返回新对象（++ Integer值会创建新对象）
					                  \item 线程安全（final字段保证可见性）
					                  \item 可以安全缓存（Integer.valueOf()使用缓存）
				                  \end{itemize}
				            \item \textbf{性能影响}：频繁修改基本类型会创建大量临时对象
			            \end{itemize}

			      \item[D] \textbf{包装类可以提高基本数据类型的操作效率 — 错误（本题答案）}
			            \begin{itemize}[label={}, leftmargin=*, nosep]
				            \item \textbf{性能对比}：
				                  \begin{itemize}[label={}, nosep]
					                  \item \textbf{内存占用}：int占用4字节，Integer对象占用16+字节（对象头+字段+填充）
					                  \item \textbf{访问开销}：int直接访问栈内存，Integer需对象引用+解引用
					                  \item \textbf{缓存机制}：-128~127范围的Integer使用缓存（IntegerCache），但也有限制
					                  \item \textbf{计算性能}：Integer需拆箱才能计算（intValue()调用）
				                  \end{itemize}
				            \item \textbf{性能测试证据}：
				                  \begin{lstlisting}[language=Java]
// 测试1：基本类型循环
long start = System.nanoTime();
for (int i = 0; i < 100_000_000; i++) {
    int sum = i + 1;
}
// 耗时：约50ms

// 测试2：包装类循环
long start = System.nanoTime();
for (Integer i = 0; i < 100_000_000; i++) {
    Integer sum = i + 1;  // 每次都拆箱+装箱
}
// 耗时：约500ms（慢10倍）

// 字节码原因：
// 1. 每次i++调用Integer.valueOf()
// 2. 每次i.intValue()拆箱
// 3. 频繁的对象创建和GC压力
				            \end{lstlisting}
				            \item \textbf{正确理解}：
				                  \begin{itemize}[label={}, nosep]
					                  \item 包装类提供"对象能力"（泛型、集合、null）
					                  \item 但不提供"性能优势"
					                  \item 反而会有额外开销
					                  \item 适用场景：需要对象特性的地方（泛型、集合、API）
					                  \item 不适用场景：数值计算密集型算法
				                  \end{itemize}
				            \item \textbf{命题人设陷阱}：利用"包装=增强"的直觉误导
			            \end{itemize}
		      \end{itemize}

		      \textbf{【应背诵知识点>}
		      \begin{itemize}[label={}, nosep]
			      \item 包装类作用：基本类型→对象，用于泛型/集合/null
			      \item 包装类不可变：final字段，修改返回新对象
			      \item 性能影响：包装类比基本类型慢2-10倍
			      \item 选项辨析：包装=对象功能，不=性能提升，反而更慢
			      \item 缓存范围：Byte/Short/Integer/Long(-128~127)
		      \end{itemize}

		      \textbf{【命题趋势与实战策略】}

		      \textbf{考场秒杀}：看到"提高效率"→警惕陷阱，包装类反而更慢
		      \textbf{高频陷阱}：
		      \begin{itemize}[label={}, nosep]
			      \item 误认为"包装=更好"
			      \item 忽略内存/CPU开销
			      \item 忽略不可变性导致的对象创建
		      \end{itemize}
		      \textbf{记忆口诀}：包装类提供对象能力，但永远不要为了性能用它
	      \end{bluebox}

	      \vspace{1em}

	\item \textbf{以下哪个是Double类的静态方法,用于将字符串转换为Double类型？}
	      \begin{enumerate}[label=\Alph*.]
		      \item parseDouble()
		      \item valueOf()
		      \item toString()
		      \item toDouble()
	      \end{enumerate}

	      \begin{bluebox}{答案与深度解析}
		      \textbf{答案：A. parseDouble()}

		      \textbf{【核心认知框架>}
		      \begin{itemize}[label={}, leftmargin=*, nosep]
			      \item \textbf{题目本质}：考查包装类静态方法的命名模式
			      \item \textbf{关键区分点}：parse→基本类型，valueOf→包装类对象
			      \item \textbf{命题人意图}：测试对Integer与Double方法一致性的理解
		      \end{itemize}

		      \textbf{【选项原理深度拆解】}

		      \begin{itemize}[leftmargin=*, nosep, label={}]
			      \item[A] \textbf{parseDouble() — 正确（返回基本类型double）}
			            \begin{itemize}[label={}, leftmargin=*, nosep]
				            \item \textbf{方法签名}：public static double parseDouble(String s) throws NumberFormatException
				            \item \textbf{返回类型}：基本类型double
				            \item \textbf{字节码}：invokestatic返回double，装箱为Double需额外步骤
				            \item \textbf{与Integer.parseInt对称}：Double.parseDouble() ↔ Integer.parseInt()
				            \item \textbf{代码示例}：
				                  \begin{lstlisting}[language=Java]
double d1 = Double.parseDouble("3.14");
System.out.println(d1);  // 3.14

double d2 = Double.parseDouble("1e10");
System.out.println(d2);  // 1.0E10

// 异常情况
double d3 = Double.parseDouble("abc");
// NumberFormatException

// 特值解析
double d4 = Double.parseDouble("Infinity");  // 无穷大
double d5 = Double.parseDouble("NaN");       // 非数字
double d6 = Double.parseDouble("-Infinity"); // 负无穷
				            \end{lstlisting}
				            \item \textbf{底层实现}：使用FloatingDecimal.parseDouble()
			            \end{itemize}

			      \item[B] \textbf{valueOf() — 错误（返回包装类Double）}
			            \begin{itemize}[label={}, leftmargin=*, nosep]
				            \item \textbf{方法签名}：public static Double valueOf(String s)
				            \item \textbf{返回类型}：Double对象
				            \item \textbf{使用场景}：需要对象时（如集合）
				            \item \textbf{代码示例}：
				                  \begin{lstlisting}[language=Java]
Double obj1 = Double.valueOf("3.14");  // 返回Double对象
double d1 = obj1.doubleValue();        // 拆箱获取double

// 对比parseDouble
double d2 = Double.parseDouble("3.14"); // 直接返回double
				            \end{lstlisting}
			            \end{itemize}

			      \item[C] \textbf{toString() — 错误（反向转换）}
			            \begin{itemize}[label={}, leftmargin=*, nosep]
				            \item 将double转为String
				            \item Double.toString(3.14)返回"3.14"
			            \end{itemize}

			      \item[D] \textbf{toDouble() — 错误（不存在）}
			            \begin{itemize}[label={}, leftmargin=*, nosep]
				            \item Double类没有toDouble()方法
				            \item 可能与intValue()混淆（用于拆箱）
			            \end{itemize}
		      \end{itemize}

		      \textbf{【应背诵知识点>}
		      \begin{itemize}[label={}, nosep]
			      \item parseXxx(String) → 返回基本类型
			      \item valueOf(String) → 返回包装类对象
			      \item 所有包装类（除Boolean外）都有 parseXxx 方法
		      \end{itemize}

		      \textbf{【命题趋势与实战策略】}

		      \textbf{考场秒杀}：字符串转基本类型→parseXxx()
		      \textbf{对仗记忆}：parseInt/parseDouble/parseLong模式一致
	      \end{bluebox}

	      \vspace{1em}

	\item \textbf{以下关于Java中的枚举类型的说法错误的是}
	      \begin{enumerate}[label=\Alph*.]
		      \item 枚举类型是一种特殊的类
		      \item 枚举类型的成员是常量
		      \item 枚举类型可以有构造方法
		      \item 枚举类型不能实现接口
	      \end{enumerate}

	      \begin{bluebox}{答案与深度解析}
		      \textbf{答案：D. 枚举类型不能实现接口}

		      \textbf{【核心认知框架>}
		      \begin{itemize}[label={}, leftmargin=*, nosep]
			      \item \textbf{题目本质}：考查enum的本质——特殊类
			      \item \textbf{关键区分点}：enum是class，可以使用类的大部分特性
			      \item \textbf{命题人意图}：纠正"枚举=数据常量"的狭隘理解
		      \end{itemize}

		      \textbf{【选项原理深度拆解】}

		      \begin{itemize}[leftmargin=*, nosep, label={}]
			      \item[A] \textbf{枚举类型是一种特殊的类 — 正确}
			            \begin{itemize}[label={}, leftmargin=*, nosep]
				            \item \textbf{字节码证据}：
				                  \begin{lstlisting}[language=Java]
enum Color {
    RED, GREEN, BLUE
}

// javap -c Color.class 结果：
public final class Color extends java.lang.Enum<Color> {
  public static final Color RED;
  public static final Color GREEN;
  public static final Color BLUE;
  // ...
}
				            \end{lstlisting}
				            \item \textbf{继承关系}：所有枚举隐式继承java.lang.Enum<E>（final类）
				            \item \textbf{类特性支持}：
				                  \begin{itemize}[label={}, nosep]
					                  \item 可以有构造方法
					                  \item 可以有字段和方法
					                  \item 可以实现接口
					                  \item 不能继承其他类（因为已继承Enum）
				                  \end{itemize}
				            \item \textbf{编译器转换}：
				                  \begin{lstlisting}[language=Java]
// 源代码
enum Weekday { MON, TUE, WED }

// 编译器等价于：
public final class Weekday extends Enum<Weekday> {
    public static final Weekday MON = new Weekday("MON", 0);
    public static final Weekday TUE = new Weekday("TUE", 1);
    public static final Weekday WED = new Weekday("WED", 2);
    private static final Weekday[] $VALUES = {MON, TUE, WED};
}
				            \end{lstlisting}
			            \end{itemize}

			      \item[B] \textbf{枚举类型的成员是常量 — 正确}
			            \begin{itemize}[label={}, leftmargin=*, nosep]
				            \item \textbf{常量特性}：
				                  \begin{itemize}[label={}, nosep]
					                  \item public static final字段
					                  \item 所有枚举成员在类加载时初始化
					                  \item 每个枚举值是单例的（全类唯一）
				                  \end{itemize}
				            \item \textbf{线程安全}：
				                  \begin{itemize}[label={}, nosep]
					                  \item 枚举实例在类初始化时创建（类加载器保证线程安全）
					                  \item 枚举是实现线程安全单例的最佳方式
				                  \end{itemize}
				            \item \textbf{代码示例}：
				                  \begin{lstlisting}[language=Java]
enum Color {
    RED, GREEN, BLUE

    @Override
    public String toString() {
        return name().toLowerCase();
    }
}

Color c1 = Color.RED;
Color c2 = Color.RED;
System.out.println(c1 == c2);  // true (单例)
System.out.println(c1.equals(c2));  // true
System.out.println(c1);  // "red" (已覆盖toString)
				            \end{lstlisting}
			            \end{itemize}

			      \item[C] \textbf{枚举类型可以有构造方法 — 正确}
			            \begin{itemize}[label={}, leftmargin=*, nosep]
				            \item \textbf{构造方法特点}：
				                  \begin{itemize}[label={}, nosep]
					                  \item 构造方法默认private（即使不写）
					                  \item 不能是public/protected（编译错误）
					                  \item 在枚举成员声明时隐式调用
				                  \end{itemize}
				            \item \textbf{语法}：
				                  \begin{lstlisting}[language=Java]
enum Planet {
    MERCURY(3.30e23, 2.44e6),
    VENUS(4.87e24, 6.06e6),
    EARTH(5.98e24, 6.37e6);

    private final double mass;   // 质量
    private final double radius; // 半径

    // 构造方法
    private Planet(double mass, double radius) {
        this.mass = mass;
        this.radius = radius;
    }

    public double getMass() {
        return mass;
    }
}

// 使用
Planet p = Planet.EARTH;
System.out.println(p.getMass());  // 5.98E24
				            \end{lstlisting}
				            \item \textbf{字节码证据}：构造方法声明为private
			            \end{itemize}

			      \item[D] \textbf{枚举类型不能实现接口 — 错误（本题答案）}
			            \begin{itemize}[label={}, leftmargin=*, nosep]
				            \item \textbf{事实}：枚举可以实现接口（因为枚举本质上就是类）
				            \item \textbf{语法}：
				                  \begin{lstlisting}[language=Java]
interface Printable {
    void print();
}

enum Color implements Printable {
    RED {
        @Override
        public void print() {
            System.out.println("红色");
        }
    },
    GREEN {
        @Override
        public void print() {
            System.out.println("绿色");
        }
    };

    // 抽象方法，每个枚举实例必须实现
    // 或者提供默认实现
}

// 或者统一实现：
enum Size implements Printable {
    SMALL, MEDIUM, LARGE;

    @Override
    public void print() {
        System.out.println("Size: " + name());
    }
}

// 使用
Color.RED.print();  // "红色"
Size.SMALL.print(); // "Size: SMALL"
				            \end{lstlisting}
				            \item \textbf{字节码验证}：
				                  \begin{lstlisting}[language=Java]
// javap Color.class
class Color extends java.lang.Enum<Color> implements Printable {
    // ...
}
				            \end{lstlisting}
				            \item \textbf{应用场景}：
				                  \begin{itemize}[label={}, nosep]
					                  \item 为枚举添加行为
					                  \item 策略模式（每个枚举实例实现不同行为）
					                  \item 与函数式接口结合
				                  \end{itemize}
				            \item \textbf{命题陷阱}：
				                  \begin{itemize}[label={}, nosep]
					                  \item 利用"枚举不能继承"误导
					                  \item 但枚举可以实现接口（因为只能有一个父类）
					                  \item 枚举已继承Enum<E>，不能继承其他类
				                  \end{itemize}
			            \end{itemize}
		      \end{itemize}

		      \textbf{【应背诵知识点>}
		      \begin{itemize}[label={}, nosep]
			      \item 枚举本质：继承java.lang.Enum<E>的特殊类
			      \item 枚举特性：隐式final，构造器private，成员是单例常量
			      \item 枚举能力：可字段/方法/实现接口，不能继承类
			      \item 选项辨析：枚举是类，可构造，可实现接口
			      \item 命题陷阱：不能继承≠不能实现接口
		      \end{itemize}

		      \textbf{【命题趋势与实战策略】}

		      \textbf{考场秒杀}：看到"不能实现接口"→选错误（枚举可以）
		      \textbf{高频陷阱}：
		      \begin{itemize}[label={}, nosep]
			      \item 枚举能实现接口，但不能继承类
			      \item 枚举构造器必须是private
			      \item 枚举实例是单例常量
		      \end{itemize}
		      \textbf{记忆口诀}：枚举就像类，能接口不能承，构造器必私有，成员是常量
	      \end{bluebox}

	      \vspace{1em}

	\item \textbf{以下哪个类用于获取系统属性？}
	      \begin{enumerate}[label=\Alph*.]
		      \item System
		      \item Properties
		      \item Runtime
		      \item SystemProperties
	      \end{enumerate}

	      \begin{bluebox}{答案与深度解析}
		      \textbf{答案：B. Properties}

		      \textbf{【核心认知框架>}
		      \begin{itemize}[label={}, leftmargin=*, nosep]
			      \item \textbf{题目本质}：考查系统属性API的返回类型
			      \item \textbf{关键区分点}：getProperties()返回Properties对象，不是直接类
			      \item \textbf{命题人意图}：测试API的使用模式和返回类型理解
		      \end{itemize}

		      \textbf{【选项原理深度拆解】}

		      \begin{itemize}[leftmargin=*, nosep, label={}]
			      \item[A] \textbf{System — 错误（使用它的方法）}
			            \begin{itemize}[label={}, leftmargin=*, nosep]
				            \item \textbf{正确用法}：System.getProperties()返回Properties
				            \item \textbf{System类功能}：
				                  \begin{itemize}[label={}, nosep]
					                  \item 标准输入/输出/错误（in/out/err）
					                  \item 系统属性（getProperties/getProperty）
					                  \item 内存管理（gc/runFinalization）
					                  \item 时间获取（currentTimeMillis/nanoTime）
					                  \item 程序退出（exit）
				                  \end{itemize}
				            \item \textbf{类定义}：public final class System，构造器private
				            \item \textbf{代码示例}：
				                  \begin{lstlisting}[language=Java]
// 正确调用
Properties props = System.getProperties();

// 直接获取单个属性
String javaHome = System.getProperty("java.home");
String osName = System.getProperty("os.name");

// 获取所有属性
Properties allProps = System.getProperties();
Enumeration<?> keys = allProps.propertyNames();
while (keys.hasMoreElements()) {
    String key = (String) keys.nextElement();
    System.out.println(key + " = " + allProps.getProperty(key));
}
				            \end{lstlisting}
			            \end{itemize}

			      \item[B] \textbf{Properties — 正确（返回类型与存储类）}
			            \begin{itemize}[label={}, leftmargin=*, nosep]
				            \item \textbf{类继承关系}：Properties extends Hashtable<Object,Object>
				            \item \textbf{功能特点}：
				                  \begin{itemize}[label={}, nosep]
					                  \item 键值对存储（继承Hashtable）
					                  \item 支持从文件加载/存储（load/store方法）
					                  \item 提供getProperty/setProperty便捷方法
					                  \item 键和值都是String类型
				                  \end{itemize}
				            \item \textbf{常用属性}：
				                  \begin{lstlisting}[language=Java]
Properties props = System.getProperties();

// 常用系统属性
props.getProperty("java.version");        // Java版本
props.getProperty("java.home");           // JDK安装路径
props.getProperty("os.name");             // 操作系统名称
props.getProperty("os.arch");             // 系统架构
props.getProperty("user.name");           // 用户名
props.getProperty("user.home");           // 用户主目录
props.getProperty("user.dir");            // 当前工作目录
props.getProperty("file.separator");      // 文件分隔符（Windows:\, Unix:/）
props.getProperty("path.separator");      // 路径分隔符（Windows:;, Unix::）
props.getProperty("line.separator");      // 行分隔符
props.getProperty("java.class.path");     // 类路径
				            \end{lstlisting}
				            \item \textbf{字节码}：Properties继承Hashtable的put/get方法
				            \item \textbf{存储文件格式}：
				                  \begin{lstlisting}[language=Java]
// 存储到文件
Properties props = new Properties();
props.setProperty("database.url", "jdbc:mysql://localhost:3306/mydb");
props.setProperty("database.user", "admin");
props.setProperty("database.password", "password123");

try (OutputStream out = new FileOutputStream("config.properties")) {
    props.store(out, "Database Configuration");
}

// config.properties文件内容：
# Database Configuration
# Thu Mar 05 10:30:00 CST 2026
database.password=password123
database.url=jdbc:mysql\://localhost\:3306/mydb
database.user=admin

// 从文件加载
Properties props2 = new Properties();
try (InputStream in = new FileInputStream("config.properties")) {
    props2.load(in);
}
String url = props2.getProperty("database.url");
				            \end{lstlisting}
			            \end{itemize}

			      \item[C] \textbf{Runtime — 错误（不同用途）}
			            \begin{itemize}[label={}, leftmargin=*, nosep]
				            \item \textbf{Runtime类功能}：
				                  \begin{itemize}[label={}, nosep]
					                  \item 执行系统命令（exec）
					                  \item 内存管理（totalMemory/freeMemory/maxMemory）
					                  \item 添加关闭钩子（addShutdownHook）
					                  \item 加载动态库（load/loadLibrary）
				                  \end{itemize}
				            \item \textbf{获取方式}：Runtime.getRuntime()（单例模式）
				            \item \textbf{代码示例}：
				                  \begin{lstlisting}[language=Java]
Runtime rt = Runtime.getRuntime();

// 执行系统命令
Process proc = rt.exec("notepad.exe");
Process proc2 = rt.exec("cmd /c dir");

// 内存信息
long total = rt.totalMemory();
long free = rt.freeMemory();
long max = rt.maxMemory();
long used = total - free;
System.out.println("Used: " + (used / 1024 / 1024) + " MB");
System.out.println("Free: " + (free / 1024 / 1024) + " MB");
System.out.println("Max: " + (max / 1024 / 1024) + " MB");
System.out.println("Total: " + (total / 1024 / 1024) + " MB");
				            \end{lstlisting}
			            \end{itemize}

			      \item[D] \textbf{SystemProperties — 错误（不存在此类）}
			            \begin{itemize}[label={}, leftmargin=*, nosep]
				            \item SystemProperties类不在java.lang包
				            \item 可能混淆了System.getProperties()返回的Properties对象
			            \end{itemize}
		      \end{itemize}

		      \textbf{【应背诵知识点>}
		      \begin{itemize}[label={}, nosep]
			      \item System.getProperties()返回Properties对象
			      \item Properties存储键值对，支持load/store
			      \item Runtime用于执行命令、内存管理
			      \item 选项辨析：系统属性用Properties，执行命令用Runtime
		      \end{itemize}

		      \textbf{【命题趋势与实战策略】}

		      \textbf{考场秒杀}：获取系统属性→Properties
		      \textbf{高频陷阱}：Properties是类名和返回类型
		      \textbf{记忆口诀}：System取属性得Properties，Runtime做命令
	      \end{bluebox}

	      \vspace{1em}

	\item \textbf{以下关于Java中的注解的说法错误的是}
	      \begin{enumerate}[label=\Alph*.]
		      \item 注解是一种元数据
		      \item 可以自定义注解
		      \item 注解可以影响程序的逻辑
		      \item 注解可以为代码添加额外的信息
	      \end{enumerate}

	      \begin{bluebox}{答案与深度解析}
		      \textbf{答案：C. 注解可以影响程序的逻辑}

		      \textbf{【核心认知框架>}
		      \begin{itemize}[label={}, leftmargin=*, nosep]
			      \item \textbf{题目本质}：考查注解的本质——元数据而非运行逻辑
			      \item \textbf{关键区分点}：注解=描述，不=执行
			      \item \textbf{命题人意图}：纠正"注解直接改变程序行为"的误解
		      \end{itemize}

		      \textbf{【选项原理深度拆解】}

		      \begin{itemize}[leftmargin=*, nosep, label={}]
			      \item[A] \textbf{注解是一种元数据 — 正确}
			            \begin{itemize}[label={}, leftmargin=*, nosep]
				            \item \textbf{元数据定义}：关于数据的数据
				            \item \textbf{类比}：
				                  \begin{itemize}[label={}, nosep]
					                  \item 图书分类标签：描述书的类型，不改变书的内容
					                  \item 注解：描述代码的用途/约束/配置，不改变代码执行
				                  \end{itemize}
				            \item \textbf{注解作用}：
				                  \begin{itemize}[label={}, nosep]
					                  \item 编译时检查（@Override检查方法是否真的重写）
					                  \item 编译时生成代码（Lombok @Data生成getter/setter）
					                  \item 运行时反射（@Controller被Spring扫描）
					                  \item 文档生成（@deprecated标记过时）
				                  \end{itemize}
				            \item \textbf{字节码层面}：注解存储在class文件的RuntimeVisibleAnnotations属性中
			            \end{itemize}

			      \item[B] \textbf{可以自定义注解 — 正确}
			            \begin{itemize}[label={}, leftmargin=*, nosep]
				            \item \textbf{自定义注解语法}：
				                  \begin{lstlisting}[language=Java]
// 定义注解
@Retention(RetentionPolicy.RUNTIME)
@Target(ElementType.METHOD)
public @interface MyAnnotation {
    String value() default "";
    int count() default 0;
}

// 使用注解
public class Test {
    @MyAnnotation(value = "test", count = 5)
    public void method1() {}

    // 省略默认值，使用默认
    @MyAnnotation
    public void method2() {}

    // 使用value属性，可省略属性名
    @MyAnnotation("hello")
    public void method3() {}
}
				            \end{lstlisting}
				            \item \textbf{元注解}：
				                  \begin{itemize}[label={}, nosep]
					                  \item @Retention：注解生命周期（SOURCE/CLASS/RUNTIME）
					                  \item @Target：注解应用目标（TYPE/FIELD/METHOD/PARAMETER等）
					                  \item @Documented：生成Javadoc时包含注解
					                  \item @Inherited：子类可继承父类注解
				                  \end{itemize}
				            \item \textbf{示例}：
				                  \begin{lstlisting}[language=Java]
// 自定义单元测试注解
@Retention(RetentionPolicy.RUNTIME)
@Target(ElementType.METHOD)
public @interface TestCase {
    String name() default "";
    int timeout() default 5000; // 毫秒
}

// 使用
public class CalculatorTest {
    @TestCase(name = "加法测试", timeout = 1000)
    public void testAdd() {
        // 测试逻辑
    }
}
				            \end{lstlisting}
			            \end{itemize}

			      \item[C] \textbf{注解可以影响程序的逻辑 — 错误（本题答案）}
			            \begin{itemize}[label={}, leftmargin=*, nosep]
				            \item \textbf{事实}：注解本身不改变程序逻辑
				            \item \textbf{证据1：注解被忽略}
				                  \begin{lstlisting}[language=Java]
// 编译通过，但@MyAnnotation什么都没做
@MyAnnotation
public class MyAnnotationTest {
}

// 类定义
public @interface MyAnnotation {}
				            \end{lstlisting}
				            \item \textbf{证据2：需要处理器/框架读取}
				                  \begin{lstlisting}[language=Java]
// Lombok @Data需要注解处理器才能生成代码
@Data  // 不配置处理器，编译后类没有getter/setter
public class User {
    private String name;
}

// Spring @Controller需要Spring容器扫描
@Controller  // 不在Spring容器中，不会被扫描注册
public class MyController {
}

// 需要配置：
// - Lombok：添加lombok依赖，IDE安装插件
// - Spring：配置@ComponentScan扫描包
				            \end{lstlisting}
				            \item \textbf{正确理解}：
				                  \begin{itemize}[label={}, nosep]
					                  \item 注解=信息标注
					                  \item 注解处理器/框架=根据注解执行逻辑
					                  \item 类比：标签（注解）+ 分拣器（处理器）= 分类行为
					                  \item 没有处理器的注解=没有分拣器的标签=无用
				                  \end{itemize}
				            \item \textbf{命题陷阱}：
				                  \begin{itemize}[label={}, nosep]
					                  \item 学生看到@Controller改变程序行为，以为是注解的作用
					                  \item 实际是Spring容器读取注解后采取的行动
					                  \item 注解只是"信息"，容器才是"行为执行者"
				                  \end{itemize}
			            \end{itemize}

			      \item[D] \textbf{注解可以为代码添加额外的信息 — 正确}
			            \begin{itemize}[label={}, leftmargin=*, nosep]
				            \item \textbf{信息类型}：
				                  \begin{itemize}[label={}, nosep]
					                  \item 用途说明（@Service表示这是服务层）
					                  \item 配置信息（@Autowired注入依赖）
					                  \item 约束检查（@NonNull检查非空）
					                  \item 文档说明（@Deprecated标记过时）
					                  \item 版本信息（@Version版本号）
				                  \end{itemize}
				            \item \textbf{代码示例}：
				                  \begin{lstlisting}[language=Java]
// 为代码添加信息
public class Example {
    @Deprecated  // 信息：已过时
    @SuppressWarnings("unchecked")  // 信息：抑制警告
    @Override  // 信息：重写父类方法
    public void oldMethod() {}
}

// 自定义业务信息
public enum UserStatus {
    ACTIVE, INACTIVE, DELETED
}

// 为枚举添加显示信息
public @interface DisplayName {
    String value();
}

enum Status {
    @DisplayName("激活")
    ACTIVE,

    @DisplayName("未激活")
    INACTIVE,

    @DisplayName("已删除")
    DELETED
}

// 通过反射读取注解信息
for (Status s : Status.values()) {
    DisplayName dn = s.getClass().getDeclaredField(s.name())
            .getAnnotation(DisplayName.class);
    System.out.println(s.name() + " = " + dn.value());
}
// 输出：
// ACTIVE = 激活
// INACTIVE = 未激活
// DELETED = 已删除
				            \end{lstlisting}
			            \end{itemize}
		      \end{itemize}

		      \textbf{【应背诵知识点>}
		      \begin{itemize}[label={}, nosep]
			      \item 注解本质：元数据（关于代码的描述信息）
			      \item 注解作用：标注信息，不直接改变逻辑
			      \item 注解生效：需要处理器/框架读取并执行
			      \item 核心区分：注解=标签，处理器=执行者
			      \item 陷阱清单：注解不=逻辑，但框架会根据注解执行逻辑
		      \end{itemize}

		      \textbf{【命题趋势与实战策略】}

		      \textbf{考场秒杀}：看到"影响逻辑"→选错误（注解本身不改变逻辑）
		      \textbf{高频陷阱}：
		      \begin{itemize}[label={}, nosep]
			      \item 误认为注解直接改变程序行为
			      \item 忽略处理器/框架的作用
			      \item 混淆"注解信息"与"注解处理"
		      \end{itemize}
		      \textbf{记忆口诀}：注解只是贴标签，框架看到才做事
	      \end{bluebox}

	      \vspace{1em}

\end{enumerate}

\end{document}
